\subsection{Settings and Assumptions}
\label{subsec:general-argument-bound-settings}

Before defining the estimates, we fix the coupled population-reference-empirical
dynamics and the common jump schedule used throughout the comparison.

\begin{definition}
\label{definition-ode-and-jump}
Let $U_1,\ldots,U_N$ be i.i.d. with deterministic law $\mu_0$, independently of
the population innovations, and set $X_0^i=Y_0^i=U_i$. Fix deterministic jump
times $0<t_1<\cdots<t_{\NofJ}\le T$ and amplitudes $\eta_\ell\ge0$.
Write $t_0=0$ and $t_{\NofJ+1}=T$; zero jumps are allowed.
On each nonempty interval $(t_\ell,t_{\ell+1})$, $0\le\ell\le\NofJ$, let
\[
\partial_tX_t^i=-\nabla_WF(\rho_{\mathbf X_t},X_t^i),\qquad
\partial_tY_t^i=-\nabla_WF(\mu_t,Y_t^i),\qquad
\partial_t\mu_t=\operatorname{div}(\mu_t\nabla_WF(\mu_t)).
\]
The paths are right-continuous, and $\mathbf X_{t_\ell-},\mathbf Y_{t_\ell-}$
denote their actual left limits. At a jump,
\begin{equation}
\label{eq:jump-dynamics-XY}
\mathbf X_{t_\ell}=\mathbf X_{t_\ell-}+\eta_\ell\xi_{\ell,X}^N(\mathbf X_{t_\ell-}),
\qquad
\mathbf Y_{t_\ell}=\mathbf Y_{t_\ell-}+\eta_\ell\xi_{\ell,Y}^N(\mathbf Y_{t_\ell-}),
\end{equation}
where $\xi_{\ell,Y}^N$ evaluates the population field with covariance
$K_{\mu_{t_\ell-}}$, and the empirical perturbation has the transition specified
in Assumption~\ref{assumption-common-seed}. The population law is pushed forward
by the same population field. Fresh innovations are independent across jumps;
the resulting fields can have covariances depending on earlier population
randomness. Their finite coupling is constructed in
Lemma~\ref{lem:population-first-finite-coupling}.
We assume that the population path admits the following characteristic
realization: conditional on
the population innovations, the ODE and jump maps compose into a common
measurable map \(\mathcal T_t\) (and the corresponding left-trace maps) such that
\[
Y_t^i=\mathcal T_t(U_i),\qquad
\mu_t=(\mathcal T_t)_\#\mu_0.
\]
On each ODE branch, \(\mathcal T_t\) is generated by the population
characteristic ODE; at a jump, it includes the population pushforward stated
above. This is the realization used in
\cref{lem:population-first-finite-coupling,lemma-gp-jump-implies-conditional-iid}.
\end{definition}

\subsubsection{Assumptions}

The following assumptions specify the moment representation, Gaussian
perturbations, and pathwise regularity used in the comparison. The feature
parameters and the anchor norms below are fixed independently of $N,T,p$.
The localization radius and conditional population scale will be chosen
from $N,T,p$; their dependence remains in every local comparison coefficient.

\paragraph{Weighted moments and derivative norms.}
Fix $\alpha\in(0,2]$. For Hilbert spaces $\mathsf H_m$ and features
$\phi_m:\mathbb R^d\to\mathsf H_m$, write
\begin{equation}
\label{eq:weighted-moment-space}
\begin{gathered}
 M_m(\nu):=\int\phi_m(x)\,\nu(dx)\in\mathsf H_m,\qquad M(\nu):=(M_m(\nu))_{m\ge1},\\
 a_m:=e^{-c_am^{\kappa_{\rm dec}}},\qquad
 \|h\|_{\mathsf H_a}^2:=\sum_{m\ge1}a_m^2\|h_m\|_{\mathsf H_m}^2,\qquad
 \mathsf H_a:=\{h:\|h\|_{\mathsf H_a}<\infty\},
\end{gathered}
\end{equation}
where $c_a>0$, $\kappa_{\rm dec}>1$, and each feature integral is a
Bochner integral whenever it exists. Scalar features are included by
taking $\mathsf H_m=\mathbb R$ for every $m$. For finite $K>0$ and $D\ge1$, set
\begin{equation}
\label{eq:moment-classes}
\begin{aligned}
 \mathcal P_{\alpha,K}
 &:=\left\{\nu\in\mathcal P_2(\mathbb R^d):
   \sup_{q\ge1}q^{-1/\alpha}
       \left(\int\|x\|_{\mathbb R^d}^q\,\nu(dx)\right)^{1/q}\le K\right\},\\
 \mathcal P_D&:=\{\nu\in\mathcal P_2(\mathbb R^d):
                     \nu(\{x:\|x\|_{\mathbb R^d}\le D\})=1\}.
\end{aligned}
\end{equation}
Here $K$ is a variable class scale; $K_0$ in
\cref{assumption-moment-growth} will bound unconditional particle moments.

\begin{definition}[Localized moment region]
\label{def:localized-moment-region}
For finite $K>0$ and $D\ge1$, the feasible moment region is
\begin{equation}
\label{eq:localized-feasible-moment-region}
 \mathcal V_{K,D}:=\operatorname{conv}
     \{M(\nu):\nu\in\mathcal P_{\alpha,K}\cup\mathcal P_D\}.
\end{equation}
The feature condition below makes these moment vectors elements of
$\mathsf H_a$, as shown in \cref{lem:moment-integral-formulas}(i).
The convex hull here consists of finite convex combinations.
\end{definition}

For an $r$-linear map $A:E_1\times\cdots\times E_r\to E$, use the norm
\[
 \|A\|_{E_1\times\cdots\times E_r\to E}
 :=\sup_{\|v_j\|_{E_j}\le1,\ 1\le j\le r}
                  \|A[v_1,\ldots,v_r]\|_E.
\]
All Fr\'echet derivatives of $G$ below are with respect to the
weighted Hilbert norm in \eqref{eq:weighted-moment-space}.
If $\mathbf e_m u$ inserts $u\in\mathsf H_m$ in coordinate $m$ and
zero in the other coordinates, define the coordinate derivatives by
\begin{equation}
\label{eq:moment-coordinate-derivatives}
\begin{aligned}
 \langle\partial_mG(\theta),u\rangle_{\mathsf H_m}
   &:=\fD G(\theta)[\mathbf e_m u],\\
 \langle u,\partial_{mn}^2G(\theta)v\rangle_{\mathsf H_m}
   &:=\fD^2G(\theta)[\mathbf e_m u,\mathbf e_n v].
\end{aligned}
\end{equation}
Thus $\partial_mG(\theta)\in\mathsf H_m$ and
$\partial_{mn}^2G(\theta):\mathsf H_n\to\mathsf H_m$.

\begin{assumption}[Structured moments and weighted smoothness]
\label{assumption-structured-moment-smoothness}
The following conditions hold, with fixed structural parameters
independent of $N,T,p$.
\begin{enumerate}[label=(\roman*),leftmargin=*]
\item \textbf{Coordinate feature growth.}
There are $C_\phi,c_{\rm gr}>0$ and $\kappa_{\rm gr}\ge0$ with
$\kappa_{\rm dec}>\max\{1,\kappa_{\rm gr}\}$. Set
$L_m:=C_\phi m\exp(c_{\rm gr}m^{\kappa_{\rm gr}})$.
Each $\phi_m$ is $C^2$, its third Fr\'echet derivative exists, and
\begin{equation}
\label{eq:coordinate-feature-growth}
\begin{aligned}
 \|\phi_m(x)\|_{\mathsf H_m}&\le L_m(1+\|x\|_{\mathbb R^d}^m),\\
 \|\fD^r\phi_m(x)\|_{(\mathbb R^d)^r\to\mathsf H_m}
 &\le L_m(1+\|x\|_{\mathbb R^d}^{(m-r)_+}),\qquad r=1,2,3.
\end{aligned}
\end{equation}
Here $(m-r)_+:=\max\{m-r,0\}$, and $\fD^0\phi_m=\phi_m$.
For $r\ge1$ the displayed norm is the multilinear operator norm.
Continuity of $\fD^3\phi_m$ is not required.
\item \textbf{Weighted smoothness of the moment functional.}
There is one deterministic open set $\mathcal U\subset\mathsf H_a$
and one $G\in C^3(\mathcal U;\mathbb R)$ such that
\begin{equation}
\label{eq:structured-moment-open-neighbourhood}
 \bigcup_{K>0,\ D\ge1}\mathcal V_{K,D}\subset\mathcal U,\qquad
 \sup_{\theta\in\mathcal U}
 \|\fD^3G(\theta)\|_{\mathsf H_a\times\mathsf H_a\times\mathsf H_a\to\mathbb R}
 \le C_3<\infty.
\end{equation}
The set $\mathcal U$, the functional $G$, and $C_3$ do not vary with
$K,D,N,T,p$. Convexity of $\mathcal U$ is not required:
the moment segments used below lie in $\mathcal V_{K,D}$.
\end{enumerate}
\end{assumption}

\begin{definition}[Objective induced by weighted moments]
\label{def:moment-objective}
For $\nu\in\mathcal P_2(\mathbb R^d)$ whose feature integrals exist and
whose moment vector belongs to $\mathcal U$, define
\begin{equation}
\label{eq:moment-objective}
 F(\nu):=G(M(\nu)).
\end{equation}
This domain includes every $\mathcal P_{\alpha,K}$ with finite $K$ and
every $\mathcal P_D$ by \cref{lem:moment-integral-formulas}(i) and
\eqref{eq:structured-moment-open-neighbourhood}.
\end{definition}
Local bounds on $\fD G,\fD^2G$ and their coordinate derivatives will
follow from a fixed anchor and $C_3$ in
\cref{lem:moment-integral-formulas,lem:conditional-local-coefficient-growth};
no uniform bound on these lower derivatives over all of $\mathcal U$
is imposed.

\paragraph{Jump notation.}
We use the population and empirical jump models in
\eqref{eq:population-jump} and \eqref{eq:particle-jump}.
Let $\mathcal H_{\ell-}$ be the full pre-jump history and
$\mathcal G_{\rm pop}$ the population-field sigma-field. For a law $\nu$
and a particle configuration $\mathbf Z$, set
\begin{align*}
 K_\nu(x,y)&:=\int\nabla_W^2F(\nu;x,z)\nabla_W^2F(\nu;z,y)\,\nu(dz),\\
 \mathbf M^N(\rho_{\mathbf Z},\mathbf Z)_{ij}
 &:=N^{-1/2}\nabla_W^2F(\rho_{\mathbf Z};Z^i,Z^j).
\end{align*}
\begin{assumption}[Gaussian perturbation model]
\label{assumption-common-seed}
The population and empirical jumps follow the models specified above.
At each prescribed jump, the population field is generated from a fresh
innovation independent of the initial sample and all earlier innovations.
Conditional on $\mathcal H_{\ell-}$, it is centered Gaussian with covariance
$K_{\mu_{t_\ell-}}$. For every population law $\mu_t$ and empirical law
$\rho_{\mathbf X_t},\rho_{\mathbf Y_t}$ used in the comparison,
including pre-jump laws, the Hessian kernel is symmetric under transposition
of its two spatial slots, and the operator with kernel $K_\nu$ on
$L^2(\nu;\mathbb R^d)$ is positive and trace class.
A jointly measurable version of the field is used in the population pushforward.
\end{assumption}
\paragraph{Finite representation and comparison coupling.}
Kernel symmetry gives the covariance identity
\[
 [K_{\rho_{\mathbf Z}}(Z^i,Z^j)]_{i,j}
 =\mathbf M^N(\rho_{\mathbf Z},\mathbf Z)
       \mathbf M^N(\rho_{\mathbf Z},\mathbf Z)^\top.
\]
This is the finite representation of the prescribed empirical Gaussian
transition. Lemma~\ref{lem:population-first-finite-coupling} constructs a
coordinate that is standard Gaussian conditional on $\mathcal H_{\ell-}$
and realizes this transition together with the population evaluation vector.
This is a construction, rather than an additional independence assumption
on the coordinate and the reference system. Reference conditional
independence is established under $\mathcal G_{\rm pop}$ in
Lemma~\ref{lemma-gp-jump-implies-conditional-iid}.
The covariance uses the integral Hessian term; the drift curvature uses
the full Hessian.

The following tail condition ensures that the reference and empirical particles remain in a common sub-Weibull class throughout the fixed time horizon.

\begin{assumption}[Unconditional moment growth]
\label{assumption-moment-growth}
There are deterministic $K_0>0$ and $\alpha\in(0,2]$ such that
\[
\sup_{0\le t\le T}\sup_{q\ge1}q^{-1/\alpha}
\max\{(\mathbb E\|X_t^1\|_{\mathbb R^d}^q)^{1/q},
       (\mathbb E\|Y_t^1\|_{\mathbb R^d}^q)^{1/q}\}\le K_0.
\]
\end{assumption}
\paragraph{Marginal and conditional bounds.}
The coupled construction is invariant under simultaneous particle
permutations, as proved in Step 6 of
\cref{lem:population-first-finite-coupling}; hence the same bound holds
for every particle. Bounds at left
limits follow by Fatou's lemma. The class scale $K$ in
\cref{assumption-structured-moment-smoothness} ranges over all finite
values, including $K_0$. Conditional bounds with a deterministic
enlarged scale $K_{\rm pop}$ hold on the population-measurable event of
Lemma~\ref{lem:population-moment-event}; they are not assumed almost surely
with scale $K_0$. In asymptotic statements a uniform bound on $K_0$ is
specified explicitly.

We separately assume a common deterministic Lipschitz bound on each ODE
branch. This pathwise condition is used for time-grid interpolation in
\cref{lemma-uniform-sup-gamma-bound-localized-conditional} and for the statistics in \cref{subsec:trigger-stability}; it is not a consequence
of the moment assumption. The two traces at a jump belong to different
branches, so the bound does not constrain the jump increment.

\paragraph{ODE intervals.}
For fixed $T>0$, write $I_\ell=(t_\ell,t_{\ell+1})$, with $t_0=0$
and $t_{\NofJ+1}=T$. Empty intervals, including a terminal interval
after a jump at $T$, impose no condition.
\begin{assumption}[Piecewise-Lipschitz flow]
\label{assumption-uniform-concentration-piecewise-lipschitz}
There is a deterministic constant $B\ge0$ such that
\begin{equation}
\sup_{1\le i\le N}\sup_{0\le \ell\le \NofJ}\sup_{\substack{s<t\\ s,t\in I_\ell}}
\max\Bigl\{
    \frac{\norm{Y_t^i-Y_s^i}_{\mathbb R^d}}{|t-s|},
     \frac{\norm{X_t^i-X_s^i}_{\mathbb R^d}}{|t-s|},
\Bigr\}
\le
B
\qquad\text{a.s.}
\end{equation}
\end{assumption}

The final pathwise condition rules out the degenerate case where the coupling error is carried by only a few particles.

\paragraph{Discrepancy and the anti-spikiness event.}
For
\[
    \delta_{t,i}\coloneq X_t^i-Y_t^i,
    \qquad
    r_t^2\coloneq \frac1N\sum_{i=1}^N\norm{\delta_{t,i}}_{\mathbb R^d}^2,
\]
and a constant \(A_{\mathrm{sp}}\ge 1\), define $E_{T,\mathrm{sp}}(A_{\mathrm{sp}})
    \coloneq
    \left\{
    \sup_{0\le t\le T:\,r_t>0}
    \frac{\frac1N\sum_{i=1}^N\norm{\delta_{t,i}}_{\mathbb R^d}^3}{r_t^3}
    \le
    A_{\mathrm{sp}}
    \right\}.$

\begin{assumption}[Anti-spikiness on comparison paths]
\label{assumption-anti-spikiness-event}
For a fixed $A_{\mathrm{sp}}\ge1$, the pathwise comparison is restricted
to $E_{T,\mathrm{sp}}(A_{\mathrm{sp}})$. The almost-sure anti-spikiness
condition in \cref{assumption-main-coupled-particle-regularity} is the special case where this
event has probability one; no such probability bound is asserted here.
\end{assumption}
If $r_t=0$, every discrepancy is zero. Thus, including these times,
the event is equivalently specified by
\begin{equation}
\label{eq-cubic-antispike-bound}
    \frac1N\sum_{i=1}^N\norm{X_t^i-Y_t^i}_{\mathbb R^d}^3
    \le
    A_{\mathrm{sp}}r_t^3
    \qquad(0\le t\le T).
\end{equation}

\subsubsection{Local consequences used in the comparison}

We first establish summability of the moment vectors and the integral
formulas on the regions $\mathcal V_{K,D}$ in
\cref{def:localized-moment-region}. These formulas then give one common
envelope for the localized comparison.

\paragraph{A common weighted-series estimate.}
Define
\begin{equation}
\label{eq:weighted-series-growth-function}
 \chi(u):=[\log(e+u)]^{\kappa_{\rm dec}/(\kappa_{\rm dec}-1)}\qquad(u\ge0).
\end{equation}
\begin{lemma}[Weighted-series estimate]
\label{lem:weighted-series-estimate}
With $a_m,L_m$ from \cref{assumption-structured-moment-smoothness}(i),
there is a constant $C$ depending only on the fixed feature parameters
and $\alpha$ such that, for $K>0$, $D\ge1$, and $m\ge1$,
\begin{equation}
\label{eq:weighted-series-summand-bound}
 a_mL_m\{2+D^m+K^m m^{m/\alpha}\}
 \le C\exp\{C\chi(K+D)-(c_a/4)m^{\kappa_{\rm dec}}\}.
\end{equation}
In particular,
\begin{equation}
\label{eq:weighted-series-sum-bound}
 \sum_{m\ge1}a_mL_m\{2+D^m+K^m m^{m/\alpha}\}
 \le Ce^{C\chi(K+D)}.
\end{equation}
For each fixed $c>0$ there is also a constant $C_c$ such that
\begin{equation}
\label{eq:weighted-series-scale-comparison}
 \chi(c(1+u)+1)\le C_c\chi(u)\qquad(u\ge0).
\end{equation}
\end{lemma}
\begin{proof}
Write $\kappa=\kappa_{\rm dec}$ and $s=\kappa/(\kappa-1)$ in this proof.
For $u,v\ge0$ and $\varepsilon>0$, Young's inequality with conjugate
exponents $\kappa,s$ gives
\begin{equation}
\label{eq:young-weighted-series}
 uv\le\varepsilon u^\kappa+
       \frac{\kappa-1}{\kappa}(\kappa\varepsilon)^{-1/(\kappa-1)}v^s.
\end{equation}
Indeed, substitute $a=(\kappa\varepsilon)^{1/\kappa}u$ and
$b=(\kappa\varepsilon)^{-1/\kappa}v$ into
$ab\le a^\kappa/\kappa+b^s/s$.
Taking $u=m$, $v=\log(e+K+D)$, and $\varepsilon=c_a/4$ yields
\[
 m\log(e+K+D)\le\tfrac{c_a}{4}m^\kappa+C\chi(K+D).
\]
The two remaining growth terms satisfy
\[
 c_{\rm gr}m^{\kappa_{\rm gr}}\le\tfrac{c_a}{4}m^\kappa+C,
 \qquad (m/\alpha)\log m+\log(4C_\phi m)
                  \le\tfrac{c_a}{4}m^\kappa+C.
\]
These inequalities hold for all $m\ge1$ after increasing the fixed $C$:
the ratios $m^{\kappa_{\rm gr}}/m^\kappa$ and
$m\log m/m^\kappa$ tend to zero, and the finitely many remaining values
have a finite maximum.
Since
$2+D^m+K^m m^{m/\alpha}\le4(e+K+D)^m m^{m/\alpha}$, substitution of
these three bounds gives
\[
\begin{aligned}
 &\log\!\left[a_mL_m\{2+D^m+K^m m^{m/\alpha}\}\right]\\
 &\quad\le-c_am^\kappa+c_{\rm gr}m^{\kappa_{\rm gr}}
       +m\log(e+K+D)+(m/\alpha)\log m+\log(4C_\phi m)\\
 &\quad\le(-c_a+c_a/4+c_a/4+c_a/4)m^\kappa+C\chi(K+D)+C\\
 &\quad=-(c_a/4)m^\kappa+C\chi(K+D)+C.
\end{aligned}
\]
Exponentiating and absorbing the fixed factor $e^C$ yields
\eqref{eq:weighted-series-summand-bound}.

Summing leaves $\sum_{m\ge1}e^{-(c_a/4)m^\kappa}<\infty$, proving
\eqref{eq:weighted-series-sum-bound}. Finally,
\[
 \log(e+c(1+u)+1)\le\log(c+2)+\log(e+u)
 \le(1+\log(c+2))\log(e+u).
\]
Raising to the power $s$ proves \eqref{eq:weighted-series-scale-comparison}.
The constants here do not depend on the anchor norms or on $K,D,N,T,p$.
\end{proof}

\paragraph{Fixed anchor.}
Let $\delta_0$ be the Dirac measure at the origin. Since
$\delta_0\in\mathcal P_D$, set
\begin{equation}
\label{eq:moment-anchor}
\begin{gathered}
 \theta_0:=M(\delta_0)=(\phi_m(0))_{m\ge1}\in\mathcal V_{K,D}\subset\mathcal U,\\
 g_0:=\|\fD G(\theta_0)\|_{\mathsf H_a\to\mathbb R},\qquad
 h_0:=\|\fD^2G(\theta_0)\|_{\mathsf H_a\times\mathsf H_a\to\mathbb R}.
\end{gathered}
\end{equation}
Moment summability in the next lemma justifies $\theta_0\in\mathsf H_a$.
The $C^3$ property in \cref{assumption-structured-moment-smoothness}(ii)
makes $g_0,h_0$ finite; they are fixed independently of $K,D,N,T,p$.

\begin{lemma}[Moment summability and integral formulas]
\label{lem:moment-integral-formulas}
Under \cref{assumption-structured-moment-smoothness}, for finite $K>0$
and $D\ge1$ the following statements hold.
\begin{enumerate}[label=(\roman*),leftmargin=*]
\item \textbf{Moment vectors.}
All feature integrals exist for $\nu\in\mathcal P_{\alpha,K}\cup\mathcal P_D$,
and
\begin{equation}
\label{eq:weighted-population-moment-summability}
\begin{aligned}
 \|M(\nu)\|_{\mathsf H_a}
 &\le\sum_{m\ge1}a_mL_m(1+K^m m^{m/\alpha})<\infty
       &&(\nu\in\mathcal P_{\alpha,K}),\\
 \|M(\nu)\|_{\mathsf H_a}
 &\le\sum_{m\ge1}a_mL_m(1+D^m)<\infty
       &&(\nu\in\mathcal P_D).
\end{aligned}
\end{equation}
\item \textbf{Integral formulas.}
For $\theta,\vartheta\in\mathcal V_{K,D}$, put
$\Delta:=\vartheta-\theta$ and $\theta_s:=(1-s)\theta+s\vartheta$.
Then
\begin{align}
 G(\vartheta)-G(\theta)
 &=\int_0^1\fD G(\theta_s)[\Delta]\,ds,
 \label{eq:moment-value-integral}\\
 \fD G(\vartheta)-\fD G(\theta)
 &=\int_0^1\fD^2G(\theta_s)[\Delta,\cdot]\,ds,
 \label{eq:moment-gradient-integral}\\
 \|\fD^2G(\vartheta)-\fD^2G(\theta)\|_{\mathsf H_a\times\mathsf H_a\to\mathbb R}
 &\le C_3\|\vartheta-\theta\|_{\mathsf H_a}.
 \label{eq:moment-hessian-lipschitz}
\end{align}
\item \textbf{Bounds from the anchor.}
With $\theta_0,g_0,h_0$ from \eqref{eq:moment-anchor}, for
$\theta\in\mathcal V_{K,D}$,
\begin{equation}
\label{eq:anchored-weighted-derivative-bounds}
\begin{aligned}
 \|\fD^2G(\theta)\|_{\mathsf H_a\times\mathsf H_a\to\mathbb R}
 &\le h_0+C_3\|\theta-\theta_0\|_{\mathsf H_a},\\
 \|\fD G(\theta)\|_{\mathsf H_a\to\mathbb R}
 &\le g_0+h_0\|\theta-\theta_0\|_{\mathsf H_a}
          +\tfrac12C_3\|\theta-\theta_0\|_{\mathsf H_a}^2.
\end{aligned}
\end{equation}
\end{enumerate}
\end{lemma}
\begin{proof}
\noindent\textbf{Step 1: sum the weighted coordinate moments.}
For $\nu\in\mathcal P_{\alpha,K}$, the $q=m$ case of
\eqref{eq:moment-classes} and the feature bound
\eqref{eq:coordinate-feature-growth} give
\[
 \int\|\phi_m(x)\|_{\mathsf H_m}\,\nu(dx)
 \le L_m\left(1+\int\|x\|_{\mathbb R^d}^m\,\nu(dx)\right)
 \le L_m(1+K^m m^{m/\alpha}).
\]
Thus the Bochner integrals exist and obey the same bound in norm.
For their sum, apply \cref{lem:weighted-series-estimate},
\eqref{eq:weighted-series-sum-bound}, with $D=1$ and the given $K$.
For every integer $n$,
\[
 \sum_{m\le n}a_m^2\|M_m(\nu)\|_{\mathsf H_m}^2
 \le\left(\sum_{m\le n}a_m\|M_m(\nu)\|_{\mathsf H_m}\right)^2
 \le\left(\sum_{m\ge1}a_mL_m(1+K^m m^{m/\alpha})\right)^2.
\]
Letting $n\to\infty$ proves the first line of
\eqref{eq:weighted-population-moment-summability}.
For $\nu\in\mathcal P_D$, replace the moment bound by
$\int\|x\|_{\mathbb R^d}^m\,\nu(dx)\le D^m$; apply \eqref{eq:weighted-series-sum-bound} with $K=1$ and the
given $D$ to get a finite weighted sum. The same finite-partial-sum
inequality proves the second line of \eqref{eq:weighted-population-moment-summability}.

\par\noindent\textbf{Step 2: integrate on a feasible moment segment.}
By \cref{def:localized-moment-region} and
\eqref{eq:structured-moment-open-neighbourhood},
$\theta_s\in\mathcal V_{K,D}\subset\mathcal U$ for every $s\in[0,1]$.
Since $\dot\theta_s=\Delta$, the chain rule gives
$\frac{d}{ds}G(\theta_s)=\fD G(\theta_s)[\Delta]$ and
$\frac{d}{ds}\fD G(\theta_s)=\fD^2G(\theta_s)[\Delta,\cdot]$.
Integrating gives
\cref{eq:moment-value-integral,eq:moment-gradient-integral}.
For $u,v\in\mathsf H_a$, the same calculation and assumption (ii) give
\[
\begin{aligned}
 (\fD^2G(\vartheta)-\fD^2G(\theta))[u,v]
 &=\int_0^1\fD^3G(\theta_s)[\Delta,u,v]\,ds,\\
 |(\fD^2G(\vartheta)-\fD^2G(\theta))[u,v]|
 &\le C_3\|\Delta\|_{\mathsf H_a}\|u\|_{\mathsf H_a}\|v\|_{\mathsf H_a}.
\end{aligned}
\]
The supremum over unit $u,v$ is
\eqref{eq:moment-hessian-lipschitz}.

\par\noindent\textbf{Step 3: integrate from the fixed anchor.}
Both $\theta_0$ and $\theta$ lie in $\mathcal V_{K,D}$.
Equation \eqref{eq:moment-hessian-lipschitz} first gives
$\|\fD^2G(\theta)\|_{\mathsf H_a\times\mathsf H_a\to\mathbb R}\le h_0+C_3\|\theta-\theta_0\|_{\mathsf H_a}$.
Apply this bound at $\theta_0+s(\theta-\theta_0)$ in
\eqref{eq:moment-gradient-integral}:
\[
\begin{aligned}
 \|\fD G(\theta)\|_{\mathsf H_a\to\mathbb R}
 &\le g_0+\|\theta-\theta_0\|_{\mathsf H_a}
       \int_0^1(h_0+C_3s\|\theta-\theta_0\|_{\mathsf H_a})\,ds\\
 &=g_0+h_0\|\theta-\theta_0\|_{\mathsf H_a}
          +\tfrac12 C_3\|\theta-\theta_0\|_{\mathsf H_a}^2.
\end{aligned}
\]
This proves \eqref{eq:anchored-weighted-derivative-bounds}.
\end{proof}

\paragraph{Particle localization and the conditional population scale.}
For $D\ge1$, define
\begin{equation}
\label{eq:sample-localization-event}
 E_{T,D}:=\left\{\max_{1\le i\le N}\sup_{0\le t\le T}
                    (\|X_t^i\|_{\mathbb R^d}\vee\|Y_t^i\|_{\mathbb R^d})\le D\right\}.
\end{equation}
The population quantities used next are supplied by
\cref{lem:population-moment-event}:
\begin{itemize}[leftmargin=*]
\item $K_{\rm pop}$ is the deterministic conditional moment scale in
\eqref{eq:population-exponential-scale}. It depends on the probability
budget and horizon, and need not equal the unconditional scale $K_0$.
\item $E_{\rm pop}\in\mathcal G_{\rm pop}$ is the population-measurable
event constructed there. For a budget $p_0\in(0,1)$ it has probability
at least $1-p_0$, and \eqref{eq:population-uniform-conditional-moments}
gives $\mu_t\in\mathcal P_{\alpha,K_{\rm pop}}$ on this event.
\end{itemize}
All time statements include pre-jump and post-jump values. The next
lemma records how these conclusions place the comparison paths in
the domain of $G$; the quantitative construction of $K_{\rm pop}$
and $E_{\rm pop}$ remains in \cref{lem:population-moment-event}.

\begin{lemma}[Moment points of the population, particle clouds, and interpolations]
\label{lemma-reachable-moment-tube}
With $K=K_{\rm pop}$ in \cref{def:localized-moment-region}:
\begin{enumerate}[label=(\roman*),leftmargin=*]
\item \textbf{Population.} On $E_{\rm pop}$, for every $t\le T$,
\[
 \sup_{q\ge1}q^{-1/\alpha}
       \left(\int\|x\|_{\mathbb R^d}^q\mu_t(dx)\right)^{1/q}\le K,
 \qquad \mu_t\in\mathcal P_{\alpha,K}.
\]
\item \textbf{Empirical measures and particle interpolation.}
On $E_{T,D}$, set $Z_t^{(s),i}:=(1-s)Y_t^i+sX_t^i$. Then
\[
 \rho_{\mathbf X_t},\quad \rho_{\mathbf Y_t},\quad
 \rho_{\mathbf Z_t^{(s)}}\in\mathcal P_D
 \qquad(t\le T,\ 0\le s\le1).
\]
\item \textbf{Segments in moment space.}
On $E_{\rm pop}\cap E_{T,D}$, let $\rho,\nu$ be any two measures in
\[
 \{\mu_t,\rho_{\mathbf X_t},\rho_{\mathbf Y_t}\}
       \cup\{\rho_{\mathbf Z_t^{(s)}}:0\le s\le1\}.
\]
For every $u\in[0,1]$,
\[
 (1-u)M(\rho)+uM(\nu)
 \in\mathcal V_{K,D}\subset\mathcal U.
\]
Thus the integral formulas and Hessian-difference bound in
\cref{lem:moment-integral-formulas} apply on this segment.
\end{enumerate}
\end{lemma}
\begin{proof}
For (i), use $\mu_t=\operatorname{Law}(Y_t^1\mid\mathcal G_{\rm pop})$
and the conditional moment bound
\eqref{eq:population-uniform-conditional-moments} of
\cref{lem:population-moment-event}.
For (ii), $E_{T,D}$ gives
\[
 \|Z_t^{(s),i}\|_{\mathbb R^d}\le(1-s)\|Y_t^i\|_{\mathbb R^d}+s\|X_t^i\|_{\mathbb R^d}\le D.
\]
Each stated empirical measure is therefore supported in the closed ball.
The same holds for left limits because that ball is closed.
For (iii), each endpoint belongs to \eqref{eq:localized-feasible-moment-region}
by (i)--(ii), and
\[
 (1-u)M(\rho)+uM(\nu)=M((1-u)\rho+u\nu)\in\mathcal V_{K,D}.
\]
The inclusion $\mathcal V_{K,D}\subset\mathcal U$ is exactly
\eqref{eq:structured-moment-open-neighbourhood}; no passage to a
closure is needed.
The particle interpolation in (ii) is handled by its support bound;
it is not identified with the measure mixture in the last display.
\end{proof}

\begin{definition}[One local comparison envelope]
\label{def:localized-envelope}
For $K>0$, $D\ge1$, with $a_m,L_m$ from
\cref{eq:weighted-moment-space,eq:coordinate-feature-growth} and the
anchor norms $g_0,h_0$ from \eqref{eq:moment-anchor}, define
\begin{equation}
\label{eq:conditional-comparison-envelope}
 \mathrm L_{K,D}:=4(1+g_0+h_0+C_3)
 \left[1+\sum_{m\ge1}a_mL_m
       \{2+D^m+K^m m^{m/\alpha}\}\right]^3.
\end{equation}
In the localized comparison $K=K_{\rm pop}$ is deterministic, and we write
$\mathrm L_D:=\mathrm L_{K_{\rm pop},D}$. This abbreviation retains the
dependence of $K_{\rm pop}$ on $N,T,p$; $\mathrm L_D$ is not assumed $O(1)$.
\end{definition}

\begin{lemma}[Growth and derivative bounds for the local envelope]
\label{lem:conditional-local-coefficient-growth}
Assume \cref{assumption-structured-moment-smoothness}. For $K>0$, $D\ge1$:
\begin{enumerate}[label=(\roman*),leftmargin=*]
\item \textbf{Growth in the localization scales.} There is a constant $C$
depending only on the fixed feature parameters, $\alpha$, $g_0,h_0,C_3$,
and the fixed dimension, such that
\begin{equation}
\label{eq:local-envelope-growth}
 1\le\mathrm L_{K,D}\le C\exp\{C\chi(K+D)\}.
\end{equation}
In particular $C$ is independent of $K,D,N,T,p$.
\item \textbf{Moment derivatives on the feasible region.} For
$\theta\in\mathcal V_{K,D}$,
\begin{equation}
\label{eq:local-moment-derivative-bound}
 \|\fD G(\theta)\|_{\mathsf H_a\to\mathbb R}
 \vee\|\fD^2G(\theta)\|_{\mathsf H_a\times\mathsf H_a\to\mathbb R}
 \le\mathrm L_{K,D}.
\end{equation}
\end{enumerate}
\end{lemma}
\begin{proof}
Use the following quantities only within this proof:
\[
 R:=1+\sum_{m\ge1}a_mL_m\{2+D^m+K^m m^{m/\alpha}\},
 \qquad Q:=1+g_0+h_0+C_3.
\]
Thus $\mathrm L_{K,D}=4QR^3$.

\noindent\textbf{Step 1: apply the common series estimate.}
By \cref{lem:weighted-series-estimate}, \eqref{eq:weighted-series-sum-bound},
$R\le Ce^{C\chi(K+D)}$. Substitution into
\eqref{eq:conditional-comparison-envelope} gives
$\mathrm L_{K,D}=4QR^3\le Ce^{C\chi(K+D)}$, with the fixed anchor
factor absorbed into $C$. This proves (i) without changing the defined
finite-sample coefficient $\mathrm L_{K,D}$.

\noindent\textbf{Step 2: apply the anchor estimates on the specified region.}
By \eqref{eq:coordinate-feature-growth}, $\|\phi_m(0)\|_{\mathsf H_m}\le L_m$.
For either class in \eqref{eq:moment-classes},
\[
\begin{aligned}
 \|M_m(\nu)-\phi_m(0)\|_{\mathsf H_m}
 &\le\int\|\phi_m(x)\|_{\mathsf H_m}\,\nu(dx)+\|\phi_m(0)\|_{\mathsf H_m}\\
 &\le L_m\left(2+\int\|x\|_{\mathbb R^d}^m\,\nu(dx)\right)\\
 &\le
 \begin{cases}
 L_m(2+D^m),&\nu\in\mathcal P_D,\\
 L_m(2+K^m m^{m/\alpha}),&\nu\in\mathcal P_{\alpha,K}.
 \end{cases}
\end{aligned}
\]
For $\theta\in\mathcal V_{K,D}$, the finite-convex-hull definition
\eqref{eq:localized-feasible-moment-region} gives
$\theta=\sum_{j=1}^J\lambda_jM(\nu_j)$, where
$\lambda_j\ge0$, $\sum_j\lambda_j=1$, and
$\nu_j\in\mathcal P_{\alpha,K}\cup\mathcal P_D$.
Using the weighted $\ell^2\le\ell^1$ estimate from
\cref{lem:moment-integral-formulas}(i), we get
\[
\begin{aligned}
 \|\theta-\theta_0\|_{\mathsf H_a}
 &\le\sum_{j=1}^J\lambda_j
           \sum_{m\ge1}a_m\|M_m(\nu_j)-\phi_m(0)\|_{\mathsf H_m}\\
 &\le\sum_{j=1}^J\lambda_j
           \sum_{m\ge1}a_mL_m\{2+D^m+K^m m^{m/\alpha}\}\\
 &=R-1\le R.
\end{aligned}
\]
Substitute this bound into
\eqref{eq:anchored-weighted-derivative-bounds}. Since $R\ge1$,
$h_0+C_3R\le(h_0+C_3)R\le QR$ and
$g_0+h_0R+(C_3/2)R^2\le(g_0+h_0+C_3)R^2\le QR^2$.
Thus
\begin{equation}
\label{eq:local-anchor-intermediate}
\begin{aligned}
 \sup_{\theta\in\mathcal V_{K,D}}
     \|\fD G(\theta)\|_{\mathsf H_a\to\mathbb R}&\le QR^2,\\
 \sup_{\theta\in\mathcal V_{K,D}}
     \|\fD^2G(\theta)\|_{\mathsf H_a\times\mathsf H_a\to\mathbb R}&\le QR.
\end{aligned}
\end{equation}
Both are at most $4QR^3=\mathrm L_{K,D}$, proving (ii).
To obtain coordinate decay, recall
\eqref{eq:moment-coordinate-derivatives} and
$\|\mathbf e_m u\|_{\mathsf H_a}=a_m\|u\|_{\mathsf H_m}$.
For unit $u\in\mathsf H_m$, $v\in\mathsf H_n$,
\[
\begin{aligned}
 |\langle\partial_mG(\theta),u\rangle|
 &=|\fD G(\theta)[\mathbf e_m u]|
   \le\mathrm L_{K,D}a_m,\\
 |\langle u,\partial_{mn}^2G(\theta)v\rangle|
 &=|\fD^2G(\theta)[\mathbf e_m u,\mathbf e_n v]|
   \le\mathrm L_{K,D}a_ma_n.
\end{aligned}
\]
Taking the unit suprema gives
$\|\partial_mG(\theta)\|_{\mathsf H_m}\le\mathrm L_{K,D}a_m$ and
$\|\partial_{mn}^2G(\theta)\|_{\mathsf H_n\to\mathsf H_m}
\le\mathrm L_{K,D}a_ma_n$.
\end{proof}

\begin{lemma}[Weighted feature regularity from coordinate bounds]
\label{lem:localized-weighted-feature-envelopes}
Under item (i) of \cref{assumption-structured-moment-smoothness}, the map
$J_0(x):=(\phi_m(x))_{m\ge1}$ is a $C^2$ map from $\mathbb R^d$ to
$\mathsf H_a$, with
\[
 \fD J_0(x)[h]=(\fD\phi_m(x)[h])_{m\ge1},\qquad
 \fD^2J_0(x)[h,v]=(\fD^2\phi_m(x)[h,v])_{m\ge1}.
\]
For $D\ge1$,
\begin{align}
 \sup_{\|x\|_{\mathbb R^d}\le D}\|\fD^rJ_0(x)\|_{(\mathbb R^d)^r\to\mathsf H_a}
 &\le\sum_{m\ge1}a_mL_m(1+D^m)<\infty
       \quad(r=0,1,2),\label{eq:weighted-feature-derivative-bound}\\
 \|\fD^rJ_0(x)-\fD^rJ_0(y)\|_{(\mathbb R^d)^r\to\mathsf H_a}
 &\le\left(\sum_{m\ge1}a_mL_m(1+D^m)\right)\|x-y\|_{\mathbb R^d}
       \quad(r=1,2).
 \label{eq:weighted-feature-spatial-difference}
\end{align}
The difference bound holds for $\|x\|_{\mathbb R^d},\|y\|_{\mathbb R^d}\le D$.
The $r=0$ subscript denotes the $\mathsf H_a$ norm; for $r=1,2$ it is the
multilinear operator norm from $(\mathbb R^d)^r$ to $\mathsf H_a$.
\end{lemma}
\begin{proof}
\noindent\textbf{Step 1: bound the coordinate maps and their differences.}
For fixed $D$, summability of
$\sum_m a_mL_m(1+D^m)$ is
\eqref{eq:weighted-population-moment-summability} for bounded support.
For $r=0,1,2$, \eqref{eq:coordinate-feature-growth} and $D\ge1$ give, for unit directions,
\[
 \|(\fD^r\phi_m(x)[v_1,\ldots,v_r])_m\|_{\mathsf H_a}
 \le\left(\sum_m a_m^2L_m^2(1+D^m)^2\right)^{1/2}
 \le\sum_m a_mL_m(1+D^m).
\]
For $r=0$ the left-hand side means $\|J_0(x)\|_{\mathsf H_a}$.
Thus the first- and second-derivative candidate sequences define bounded
linear and bilinear maps, respectively.
If $\|x\|_{\mathbb R^d},\|y\|_{\mathbb R^d}\le D$, then
$\|(1-s)y+sx\|_{\mathbb R^d}\le(1-s)\|y\|_{\mathbb R^d}+s\|x\|_{\mathbb R^d}\le D$.
The mean-value inequality, applied to $\fD^r\phi_m$ for $r=1,2$, gives
\begin{equation}
\label{eq:coordinate-feature-spatial-difference}
\begin{aligned}
 \|\fD^r\phi_m(x)-\fD^r\phi_m(y)\|_{(\mathbb R^d)^r\to\mathsf H_m}
 &\le\sup_{\|z\|_{\mathbb R^d}\le D}\|\fD^{r+1}\phi_m(z)\|_{(\mathbb R^d)^{r+1}\to\mathsf H_m}\,\|x-y\|_{\mathbb R^d}
 \\ &\le L_m(1+D^m)\|x-y\|_{\mathbb R^d}.
\end{aligned}
\end{equation}
The norms here are from $(\mathbb R^d)^r$ to $\mathsf H_m$.
To apply the scalar mean-value theorem, evaluate on unit directions
and pair the output with a unit vector of $\mathsf H_m$; then take the
supremum. This uses the bound in \eqref{eq:coordinate-feature-growth},
including $r=3$, but no continuity of $\fD^3\phi_m$.

\par\noindent\textbf{Step 2: sum the coordinate Taylor remainders.}
Fix $D>\max\{1,\|x\|_{\mathbb R^d}\}$ and $\|h\|_{\mathbb R^d}<D-\|x\|_{\mathbb R^d}$; the segment $x+sh$
is in the ball by the convexity calculation in Step 1.
Coordinate-wise integration gives
\[
\begin{aligned}
 \phi_m(x+h)-\phi_m(x)-\fD\phi_m(x)[h]
 &=\int_0^1\{\fD\phi_m(x+sh)-\fD\phi_m(x)\}[h]\,ds,\\
 \{\fD\phi_m(x+h)-\fD\phi_m(x)\}[v]-\fD^2\phi_m(x)[h,v]
 &=\int_0^1\{\fD^2\phi_m(x+sh)-\fD^2\phi_m(x)\}[h,v]\,ds.
\end{aligned}
\]
Apply \eqref{eq:coordinate-feature-spatial-difference} with points $x+sh,x$,
whose distance is $s\|h\|_{\mathbb R^d}$. For unit $v$, each right-hand side has norm
at most $L_m(1+D^m)\|h\|_{\mathbb R^d}^2\int_0^1s\,ds
=\tfrac12 L_m(1+D^m)\|h\|_{\mathbb R^d}^2$.
The weighted norm of a sequence is at most the weighted sum of its
coordinate norms, by the finite-partial-sum inequality
$(\sum_{m\le n}b_m^2)^{1/2}\le\sum_{m\le n}b_m$ for $b_m\ge0$
and then $n\to\infty$. Consequently,
\[
\begin{aligned}
 &\|J_0(x+h)-J_0(x)-(\fD\phi_m(x)[h])_m\|_{\mathsf H_a}\\
 &\qquad\le\tfrac12\left(\sum_m a_mL_m(1+D^m)\right)\|h\|_{\mathbb R^d}^2,\\
 &\sup_{\|v\|_{\mathbb R^d}\le1}
 \|(\{\fD\phi_m(x+h)-\fD\phi_m(x)\}[v]-\fD^2\phi_m(x)[h,v])_m\|_{\mathsf H_a}\\
 &\qquad\le\tfrac12\left(\sum_m a_mL_m(1+D^m)\right)\|h\|_{\mathbb R^d}^2.
\end{aligned}
\]
Dividing by $\|h\|_{\mathbb R^d}$ proves the two derivative identities. Summing
\eqref{eq:coordinate-feature-spatial-difference} in the same weighted
norm gives \eqref{eq:weighted-feature-spatial-difference}, so these
derivatives are continuous and $J_0$ is $C^2$.
\end{proof}

\begin{lemma}[Local comparison bounds with common moment weights]
\label{lemma-localized-moment-smoothness-consequences}
Assume \cref{assumption-structured-moment-smoothness}. Fix $K>0$, $D\ge1$,
and define
\begin{equation}
\label{eq:common-local-moment-weights}
 \gamma_m(K,D):=\mathrm L_{K,D}a_m
              =\mathrm L_{K,D}e^{-c_am^{\kappa_{\rm dec}}}.
\end{equation}
Let $\rho,\nu$ be in the domain of $F$ in
\cref{eq:moment-objective}, with
$M(\rho),M(\nu)\in\mathcal V_{K,D}$, and let $x,x',y,y'$ be in
the ball of radius $D$. Write
$\Delta_m:=\|M_m(\rho)-M_m(\nu)\|_{\mathsf H_m}$.
\begin{enumerate}[label=(\roman*),leftmargin=*]
\item \textbf{Dependence on the measure.}
\begin{align}
 |F(\rho)-F(\nu)|\vee
 \|\nabla_WF(\rho,x)-\nabla_WF(\nu,x)\|_{\mathbb R^d}
 &\le\sum_m\gamma_m(K,D)\Delta_m,
 \label{eq:local-value-gradient-moment-bound}\\
 \max\bigl\{\|\nabla_x\nabla_WF(\rho,x)-\nabla_x\nabla_WF(\nu,x)\|_{\mathbb R^d\to\mathbb R^d},
 &\notag\\[-0.5ex]
 \|\nabla_W^2F(\rho;x,y)-\nabla_W^2F(\nu;x,y)\|_{\mathbb R^d\to\mathbb R^d}\bigr\}
 &\le\tfrac12\sum_m\gamma_m(K,D)\Delta_m.
 \label{eq:local-two-hessians-moment-bound}
\end{align}
In particular the \emph{sum} of the two Hessian differences in
\eqref{eq:local-two-hessians-moment-bound} is at most
$\sum_m\gamma_m(K,D)\Delta_m$.
\item \textbf{Magnitude and spatial dependence.}
\begin{equation}
\label{eq:local-field-magnitude-bound}
 \|\nabla_WF(\nu,x)\|_{\mathbb R^d}\vee
 \|\nabla_x\nabla_WF(\nu,x)\|_{\mathbb R^d\to\mathbb R^d}\vee
 \|\nabla_W^2F(\nu;x,y)\|_{\mathbb R^d\to\mathbb R^d}\le\mathrm L_{K,D}.
\end{equation}
The spatial differences satisfy
\begin{align}
 \|\nabla_WF(\nu,x)-\nabla_WF(\nu,x')\|_{\mathbb R^d}
 &\le\mathrm L_{K,D}\|x-x'\|_{\mathbb R^d},
 \label{eq:local-gradient-spatial-bound}\\
 \|\nabla_x\nabla_WF(\nu,x)-\nabla_x\nabla_WF(\nu,x')\|_{\mathbb R^d\to\mathbb R^d}
 &\le\mathrm L_{K,D}\|x-x'\|_{\mathbb R^d},
 \label{eq:local-jacobian-spatial-bound}\\
 \|\nabla_W^2F(\nu;x,y)-\nabla_W^2F(\nu;x',y')\|_{\mathbb R^d\to\mathbb R^d}
 &\le\mathrm L_{K,D}(\|x-x'\|_{\mathbb R^d}+\|y-y'\|_{\mathbb R^d}).
 \label{eq:local-kernel-spatial-bound}
\end{align}
\end{enumerate}
For the localized comparison we abbreviate
$\gamma_m(D):=\gamma_m(K_{\rm pop},D)=\mathrm L_Da_m$.
\end{lemma}

\begin{proof}
\noindent\textbf{Step 1: recall the derivative representations and their bounds.}
Within this proof set
\[
 Q:=1+g_0+h_0+C_3,\qquad
 R:=1+\sum_m a_mL_m\{2+D^m+K^m m^{m/\alpha}\},
 \qquad \mathrm L_{K,D}=4QR^3.
\]
On the stated moment and spatial regions,
\cref{eq:local-anchor-intermediate,eq:weighted-feature-derivative-bound} give
\begin{equation}
\label{eq:local-proof-bounds}
\begin{gathered}
 \|\fD G\|_{\mathsf H_a\to\mathbb R}\le QR^2,\qquad
 \|\fD^2G\|_{\mathsf H_a\times\mathsf H_a\to\mathbb R}\le QR,\\
 \|\fD^rJ_0(x)\|_{(\mathbb R^d)^r\to\mathsf H_a}\le R
       \qquad(r=1,2).
\end{gathered}
\end{equation}
For $\theta_\nu:=M(\nu)$ and $u,v\in\mathbb R^d$, the moment chain rule for
\cref{eq:moment-objective} is
\begin{align}
 \langle\nabla_WF(\nu,x),v\rangle
 &=\fD G(\theta_\nu)[\fD J_0(x)v],
 \label{eq:weighted-gradient-representation}\\
 \langle u,\nabla_x\nabla_WF(\nu,x)v\rangle
 &=\fD G(\theta_\nu)[\fD^2J_0(x)[u,v]],
 \label{eq:weighted-spatial-hessian-representation}\\
 \langle u,\nabla_W^2F(\nu;x,y)v\rangle
 &=\fD^2G(\theta_\nu)[\fD J_0(x)u,\fD J_0(y)v].
 \label{eq:weighted-integral-hessian-representation}
\end{align}
These are the coordinate chain-rule sums, with the directions collected
in $\mathsf H_a$. For example the first sum is
$\sum_m\langle\partial_mG(\theta_\nu),\fD\phi_m(x)v\rangle$.
Finite coordinate sums converge to the displayed forms because the
directions converge in $\mathsf H_a$ by
\cref{lem:localized-weighted-feature-envelopes}, and $\fD G,\fD^2G$ are
bounded forms on $\mathsf H_a$ by
\cref{assumption-structured-moment-smoothness}(ii).

\noindent\textbf{Step 2: change the measure along its moment segment.}
Set $\theta_\rho:=M(\rho)$, $\Delta:=\theta_\rho-\theta_\nu$, and
$\theta_s:=\theta_\nu+s\Delta\in\mathcal V_{K,D}$. By
\cref{eq:moment-value-integral,eq:moment-gradient-integral,eq:local-proof-bounds},
\begin{align*}
 |F(\rho)-F(\nu)|
 &\le\int_0^1QR^2\|\Delta\|_{\mathsf H_a}\,ds
  =QR^2\|\Delta\|_{\mathsf H_a},\\
 \|\fD G(\theta_\rho)-\fD G(\theta_\nu)\|_{\mathsf H_a\to\mathbb R}
 &\le\int_0^1QR\|\Delta\|_{\mathsf H_a}\,ds
  =QR\|\Delta\|_{\mathsf H_a}.
\end{align*}
Substituting the second inequality in
\cref{eq:weighted-gradient-representation,eq:weighted-spatial-hessian-representation}
and taking suprema over unit directions gives
\[
\begin{aligned}
 \|\nabla_WF(\rho,x)-\nabla_WF(\nu,x)\|_{\mathbb R^d}
 &\le QR\|\Delta\|_{\mathsf H_a}\|\fD J_0(x)\|_{\mathbb R^d\to\mathsf H_a}
  \le QR^2\|\Delta\|_{\mathsf H_a},\\
 \|\nabla_x\nabla_WF(\rho,x)-\nabla_x\nabla_WF(\nu,x)\|_{\mathbb R^d\to\mathbb R^d}
 &\le QR\|\Delta\|_{\mathsf H_a}\|\fD^2J_0(x)\|_{\mathbb R^d\times\mathbb R^d\to\mathsf H_a}
  \le QR^2\|\Delta\|_{\mathsf H_a}.
\end{aligned}
\]
For the integral Hessian, use
\cref{eq:moment-hessian-lipschitz,eq:weighted-integral-hessian-representation}:
\[
\begin{aligned}
 &\|\nabla_W^2F(\rho;x,y)-\nabla_W^2F(\nu;x,y)\|_{\mathbb R^d\to\mathbb R^d}\\
 &\quad\le C_3\|\Delta\|_{\mathsf H_a}
       \|\fD J_0(x)\|_{\mathbb R^d\to\mathsf H_a}
       \|\fD J_0(y)\|_{\mathbb R^d\to\mathsf H_a}
 \le QR^2\|\Delta\|_{\mathsf H_a}.
\end{aligned}
\]
Finally,
\[
 \|\Delta\|_{\mathsf H_a}\le\sum_m a_m\Delta_m,
 \qquad QR^2\le\tfrac12\mathrm L_{K,D}
 \quad\text{because }\mathrm L_{K,D}=4QR^3,\ R\ge1.
\]
This proves both inequalities in (i), with room to sum the two Hessian errors.

\noindent\textbf{Step 3: change the spatial variables.}
The three representations and \eqref{eq:local-proof-bounds} give
\[
\begin{aligned}
 \|\nabla_WF(\nu,x)\|_{\mathbb R^d}&\le QR^2R,\\
 \|\nabla_x\nabla_WF(\nu,x)\|_{\mathbb R^d\to\mathbb R^d}&\le QR^2R,\\
 \|\nabla_W^2F(\nu;x,y)\|_{\mathbb R^d\to\mathbb R^d}&\le QRR^2.
\end{aligned}
\]
All three are at most $QR^3\le\mathrm L_{K,D}$.
Equation \eqref{eq:weighted-feature-spatial-difference} gives
$\|\fD^rJ_0(x)-\fD^rJ_0(x')\|_{(\mathbb R^d)^r\to\mathsf H_a}\le R\|x-x'\|_{\mathbb R^d}$ for $r=1,2$.
Inserting these bounds in the first two representations gives
\cref{eq:local-gradient-spatial-bound,eq:local-jacobian-spatial-bound},
since $QR^2R\le\mathrm L_{K,D}$.
For the kernel, change one slot at a time:
\[
\begin{aligned}
 &\fD^2G(\theta_\nu)[\fD J_0(x)u,\fD J_0(y)v]
      -\fD^2G(\theta_\nu)[\fD J_0(x')u,\fD J_0(y')v]\\
 &=\fD^2G(\theta_\nu)[(\fD J_0(x)-\fD J_0(x'))u,\fD J_0(y)v]\\
 &\quad+\fD^2G(\theta_\nu)[\fD J_0(x')u,(\fD J_0(y)-\fD J_0(y'))v].
\end{aligned}
\]
For unit $u,v$, its absolute value is at most
$QR\{R^2\|x-x'\|_{\mathbb R^d}+R^2\|y-y'\|_{\mathbb R^d}\}$.
Taking the supremum proves \eqref{eq:local-kernel-spatial-bound}.
\end{proof}


In the local comparison estimates below we work on $E_{\rm pop}$ and use the deterministic conditional moment scale and coefficient convention from the assumptions. Spatial localization and anti-spikiness are imposed where stated.

\subsection{Bound $r_t = \frac{1}{\sqrt{N}}\norm{\bs{X}_t-\bs{Y}_t}_{2}$ between jumps}
\label{subsec:general-argument-bound-r_t-ode-continuous}

Throughout this subsection, use $\mathrm L_D=\mathrm L_{K_{\rm pop},D}$
and $\gamma_m(D)=\mathrm L_Da_m$ from
\cref{lemma-localized-moment-smoothness-consequences}. We suppress the
argument $D$ of $\gamma_m$ when the localization radius is fixed.

On each ODE interval, the discrepancy between \(\bs X_t\) and \(\bs Y_t\) is read through the segment joining the two configurations.

For $s\in[0,1]$, define the interpolating particle configuration
\begin{equation}
    \bs Z_t^{(s)} \coloneq (1-s)\bs Y_t + s\bs X_t \in \R^{dN},\quad (Z_t^{(s)})^i \coloneq (1-s)Y_t^i + sX_t^i 
\end{equation}

The linearization of the empirical vector field produces the following block Hessian matrix on \((\mathbb R^d)^N\).

For $\mu \in \mathcal{P}(\R^d)$ and $\bs{v} \in (\R^{d})^{N}$, define a block matrix
\begin{equation}
    \HessianMatrices(\mu,\bs{v})\coloneq
    \left\{
    \frac{1}{N}
    \nabla_W^2 F\!\left(\mu,v_i,v_j\right)
    + \delta_{ij}\nabla_x\nabla_W F\!\left(\mu,v_i\right)\in \R^{d\times d}\right\}_{i,j=1}^N.
\end{equation}
For each $(t,s)$, define a $\R^{dN} \times \R^{dN}$ block matrix
\begin{align}
    \HessianMatrices_t^{(s)}
    \coloneq&
    \HessianMatrices(\rho_{\bs Z_t^{(s)}},\bs{Z}_t^{(s)}).
\end{align}
For each $t\ge 0$, define the population-sampled Hessian matrix along the coupled
infinite-particle trajectory by
\begin{equation}
    \HessianMatrices_t^{Y}
    \coloneq
    \HessianMatrices(\mu_t,\bs{Y}_t).
\end{equation}

The smallest eigenvalue of the population-sampled block Hessian is the local curvature quantity that drives contraction or expansion of the coupling error.

Let
\begin{equation}
    \kappa_t^{Y,N}
    \coloneq
    \lambda_{\min}\!\bigl(\HessianMatrices_t^{Y}\bigr).
\end{equation}

\paragraph{Bound the deterministic dynamics}

The first lemma identifies the exact Jacobian of the empirical vector field, so that the particlewise difference can be expressed through a block Hessian.

\begin{lemma}[Mean-value theorem for the empirical vector field]
\label{lemma-mean-value-theorem}
Let $V_N(\bs X)
:=
\bigl(V(\rho_{\bs X},X_i)\bigr)_{i=1}^N,\;
\rho_{\bs X}
:=
\frac1N\sum_{k=1}^N \delta_{X_k}$,
where \(V:\mathcal P_2(\R^d)\times \R^d\to\R^d\) is \(C^1\) in the spatial variable and Wasserstein differentiable in the measure variable, with continuous derivatives \(\nabla_x V\) and \(\nabla_W V\).
For \(\bs X_t,\bs Y_t\in(\R^d)^N\), define
$\Delta_t:=\bs X_t-\bs Y_t,\;
\bs Z_t^{(s)}:=\bs Y_t+s\Delta_t.
$
Then
\[
V_N(\bs X_t)-V_N(\bs Y_t)
=
\int_0^1 DV_N(\bs Z_t^{(s)})\,\Delta_t\,\dd s,
\]
where \(DV_N(\bs Z)\in\R^{Nd\times Nd}\) is the block matrix whose \((i,j)\)-block is given by
\[
\bigl[DV_N(\bs Z)\bigr]_{ij}
=
\delta_{ij}\,\nabla_x V(\rho_{\bs Z},Z_i)
+
\frac1N\,\nabla_W V(\rho_{\bs Z},Z_i)(Z_j).
\]
For the moment objective \(F\) in \cref{def:moment-objective}, the same
mean-value identity holds for
\(V_N(\bs X)=(-\nabla_WF(\rho_{\bs X},X_i))_{i=1}^N\), with $\bigl[DV_N(\bs Z)\bigr]_{ij}
=
-\delta_{ij}\,\nabla_x\nabla_W F(\rho_{\bs Z},Z_i)
-\frac1N\,\nabla_W^2F(\rho_{\bs Z},Z_i,Z_j).$
\end{lemma}
\begin{proof}
\noindent\textbf{Step 1: integrate the configuration derivative.}
Fix \(t\) and define $ \Phi(s):=V_N(\bs Z_t^{(s)})\in(\R^d)^N.$
Since \(V_N\) is \(C^1\) on \((\R^d)^N\), the fundamental theorem of calculus gives
\[
V_N(\bs X_t)-V_N(\bs Y_t)
=
\Phi(1)-\Phi(0)
=
\int_0^1 \Phi'(s)\,\dd s
=
\int_0^1 DV_N(\bs Z_t^{(s)})\,\Delta_t\,\dd s.
\]
\par\noindent\textbf{Step 2: compute the spatial and measure derivatives.}
It remains to identify the Jacobian \(DV_N\).
Assume that $V(\mu,x)$ is $C^1$ in $x$ and differentiable in $\mu$ in the Wasserstein sense, with derivative denoted by $\nabla_W V(\mu,x)(\cdot)$.
Then, by the chain rule for empirical measures,
\[
\partial_{X_j}\bigl[V(\rho_{\bs X},x)\bigr]
=
\frac1N \nabla_W V(\rho_{\bs X},x)(X_j).
\]
Applying the usual chain rule to $V(\rho_{\bs X},X_i)$ yields
\[
\bigl[DV_N(\bs X)\bigr]_{ij} = \partial_{X_j}V(\rho_{\bs X},X_i)
=
\delta_{ij}\,\nabla_x V(\rho_{\bs X},X_i)
+
\frac1N\,\nabla_W V(\rho_{\bs X},X_i)(X_j).
\]
Finally, if \(V=-\nabla_W F\), then $\nabla_x V=-\nabla_x\nabla_W F,
\;
\nabla_W V=-\nabla_W^2F,$
which yields the stated block formula.
For the moment objective in \cref{def:moment-objective}, every empirical
law lies in its domain. The identity
\(M(\rho_{\bs X})=N^{-1}\sum_k J_0(X_k)\),
\cref{lem:localized-weighted-feature-envelopes}, and
\eqref{eq:weighted-gradient-representation} give a \(C^1\) empirical
field \(V_N\). Its block derivative follows directly as below. Moving only
particle $j$ in direction $h$ gives
\[
 \left.\frac d{d\varepsilon}\right|_0
 M_m\left(\frac1N\sum_{k\ne j}\delta_{X_k}
              +\frac1N\delta_{X_j+\varepsilon h}\right)
 =\frac1N\fD\phi_m(X_j)h.
\]
In the weighted space, the corresponding derivative of $M(\rho_{\bs X})$
is $N^{-1}\fD J_0(X_j)h$ by
\cref{lem:localized-weighted-feature-envelopes}. Keep the evaluation point
$x$ fixed and pair with $v\in\mathbb R^d$ in
\eqref{eq:weighted-gradient-representation}:
\[
\begin{aligned}
 \partial_{X_j}\langle\nabla_WF(\rho_{\bs X},x),v\rangle_{\mathbb R^d}[h]
 &=\fD^2G(M(\rho_{\bs X}))[N^{-1}\fD J_0(X_j)h,\fD J_0(x)v]\\
 &=N^{-1}\langle v,\nabla_W^2F(\rho_{\bs X};x,X_j)h\rangle_{\mathbb R^d}.
\end{aligned}
\]
The last equality uses symmetry of $\fD^2G$ and
\eqref{eq:weighted-integral-hessian-representation}. Since $v$ is arbitrary, set $x=X_i$ to obtain
the measure contribution $N^{-1}\nabla_W^2F(\rho_{\bs X};X_i,X_j)h$.
Moving the evaluation point contributes
$\delta_{ij}\nabla_x\nabla_WF(\rho_{\bs X},X_i)h$.
Both contributions change sign when $V=-\nabla_WF$.
\end{proof}
Using this linearization, the next lemma isolates the three forces in the ODE error: sampled curvature, Hessian mismatch, and empirical drift residual.

\begin{lemma}[Upper bound of $\dot{r}_t$ under the ODE flow given error terms $\iota_t^N$ and $\varepsilon_t^N$.]
\label{lemma-bound-r_t-given-eta-eps}
Let $V(\mu,x)\coloneq -\nabla_W F(\mu,x)$.
Consider the coupled systems
\begin{equation}
    \frac{\dd}{\dd t}X_t^i = V(\rho_{\bs X_t},X_t^i),
    \qquad
    \frac{\dd}{\dd t}Y_t^i = V(\mu_t,Y_t^i).
\end{equation}
Define
\begin{align}
    &\Delta_t \coloneq \bs X_t-\bs Y_t \in (\R^{d})^N,
    \qquad
    r_t \coloneq \frac{1}{\sqrt N}\norm{\Delta_t}_2,\\
    &\label{eq:ode-hessian-discrepancy-definition}
    \iota_t^N \coloneq \frac{1}{ N} \sup_{s\in [0,1]} |\Delta_t^\top (\HessianMatrices_t^{(s)} - \HessianMatrices_t^Y) \Delta_t|,\quad 
    \varepsilon_t^N
    \coloneq
    \left(
        \frac{1}{N}\sum_{i=1}^N
        \bigl|V(\rho_{\bs Y_t},Y_t^i)-V(\mu_t,Y_t^i)\bigr|^2
    \right)^{1/2}.
\end{align}
Assume that the map
\(
\bs x\mapsto (V(\rho_{\bs x},x^i))_{i=1}^N
\)
is differentiable along the segment joining $\bs X_t$ and $\bs Y_t$.
Then, at non-jump times with $r_t>0$,
\begin{equation}
    \dot r_t
    \le
    - \kappa_t^{Y,N} r_t
    + r_t^{-1} \iota_t^N 
    + \varepsilon_t^N.
\end{equation}
\end{lemma}

\begin{proof}
\noindent\textbf{Step 1: differentiate and split the squared discrepancy.}

First,
\begin{equation}
    \frac{1}{2}\frac{\dd}{\dd t}r_t^2
    =
    \frac{1}{N}
    \innerProd{\Delta_t}{V_N(\bs X_t)-V_\infty(\bs Y_t)},
\end{equation}
where
\begin{equation}
    V_N(\bs X_t)
    \coloneq
    \bigl(V(\rho_{\bs X_t},X_t^i)\bigr)_{i=1}^N,
    \qquad
    V_\infty(\bs Y_t)
    \coloneq
    \bigl(V(\mu_t,Y_t^i)\bigr)_{i=1}^N.
\end{equation}
We decompose
\begin{align}
    \frac{1}{2}\frac{\dd}{\dd t}r_t^2
    =&
    \underbrace{
        \frac{1}{N}
        \innerProd{\Delta_t}{V_N(\bs X_t)-V_N(\bs Y_t)}
    }_{\mathrm{(A)}}
    +
    \underbrace{
        \frac{1}{N}
        \innerProd{\Delta_t}{V_N(\bs Y_t)-V_\infty(\bs Y_t)}
    }_{\mathrm{(B)}}.
\end{align}


\par\noindent\textbf{Step 2: isolate sampled curvature and the Hessian error.}
For \(\mathrm{(A)}\), by the mean-value theorem (\cref{lemma-mean-value-theorem}),
\begin{equation}
    V_N(\bs X_t)-V_N(\bs Y_t)
    =
    -\int_0^1 \HessianMatrices_t^{(s)} \Delta_t\,\dd s.
\end{equation}
Hence
\begin{align}
    \mathrm{(A)}
    &=
    -\int_0^1
    \frac{1}{N}\Delta_t^\top \HessianMatrices_t^{(s)}\Delta_t\,\dd s \\
    &=
    -\int_0^1
    \frac{1}{N}\Delta_t^\top \HessianMatrices_t^Y \Delta_t\,\dd s
    -
    \int_0^1
    \frac{1}{N}\Delta_t^\top
    \bigl(\HessianMatrices_t^{(s)}-\HessianMatrices_t^Y\bigr)
    \Delta_t\,\dd s.
\end{align}
By the definition of \(\kappa_t\),
\begin{equation}
 \begin{aligned}
    -\frac{1}{N}\Delta_t^\top \HessianMatrices_t^{Y}\Delta_t
    &=-\innerProd{\frac{1}{\sqrt{N}}\Delta_t}{\HessianMatrices_t^Y \frac{1}{\sqrt{N}}\Delta_t}_2\\
    &\le-\kappa_t^{Y,N} \norm{\frac{1}{\sqrt{N}}\Delta_t}_{\mathbb R^{dN}}^2
     =-\kappa_t^{Y,N}r_t^2\qquad(\forall s\in[0,1]),
 \end{aligned}
\end{equation}
and                                     
\begin{equation}
    \left|\frac{1}{N}\Delta_t^\top
    \bigl(\HessianMatrices_t^{(s)}- \HessianMatrices_t^{Y}\bigr)\Delta_t
    \right|
    \le
    \iota_t^N.
\end{equation}
Therefore,
\begin{equation}
    \mathrm{(A)}
    \le
    -\kappa_t^{Y,N} r_t^2 
    + \iota_t^N.
\end{equation}


\par\noindent\textbf{Step 3: bound the drift residual and divide for positive error.}
For \(\mathrm{(B)}\), by the Cauchy--Schwarz inequality,
\begin{align}
    \mathrm{(B)}
    &\le
    \left(
        \frac{1}{N}\norm{\Delta_t}_2^2
    \right)^{1/2}
    \left(
        \frac{1}{N}
        \sum_{i=1}^N
        \bigl|
            V(\rho_{\bs Y_t},Y_t^i)-V(\mu_t,Y_t^i)
        \bigr|^2
    \right)^{1/2} \\
    &=
    r_t\varepsilon_t^N.
\end{align}

Combining the estimates for \(\mathrm{(A)}\) and \(\mathrm{(B)}\), we obtain
\begin{equation}
    \frac{1}{2}\frac{\dd}{\dd t}r_t^2
    \le
    -\kappa_t^{Y,N} r_t^2 
    + \iota_t^N
    + r_t\varepsilon_t^N.
\end{equation}
Hence, whenever \(r_t>0\),
\begin{equation}
    \dot r_t
    \le
    -\kappa_t^{Y,N} r_t 
    + \iota_t^N r_t^{-1}
    + \varepsilon_t^N.
\end{equation}
At $r_t=0$ the quotient is not used; the norm derivative is treated directly in Lemma~\ref{lemma-common-empirical-moment-residual-directional-rt}.
\end{proof}

\paragraph{Bound the error terms $\iota_t^N$ and $\varepsilon_t^N$ by moment errors.}

We will repeatedly convert block-operator estimates into scalar matrix estimates by bounding each block norm separately.

\begin{lemma}[operator norm bound of blocked matrices.]
\label{lemma-operator-norm-block-matrix}
    Let 
    \begin{equation}
        A = \begin{bmatrix}
            B_{1,1} & \dots& B_{1,N}\\
            \vdots &\ddots & \vdots \\
            B_{N,1}& \dots &B_{N,N}
        \end{bmatrix},
        B_{i,j} \in \R^{d\times d}
    \end{equation}
    be a $\R^{dN}\times \R^{dN}$ block matrix and
    \begin{equation}
        C = (\norm{B_{i,j}}_{\mathbb R^d\to\mathbb R^d})_{i,j=1}^N
    \end{equation}
    be a corresponding $\R^{N}\times \R^{N}$ matrix.
    then,
    \begin{equation}
        \norm{A}_{\mathbb R^{dN}\to\mathbb R^{dN}} \leq \norm{C}_{\mathbb R^{N}\to\mathbb R^{N}}
    \end{equation}
\end{lemma}
\begin{proof}
    For any vector $x\in \R^{dN}$ such that $\norm{x}_2=1$, we write $x = (x_1,\dots,x_N)$, $x_i\in \R^d$.
    First note that
    \begin{equation}
        1 = \norm{x}_2^2 = \sum_i (\sum_j x^2_{i,j}) = \sum_i z_i^2 = \norm{z}_2^2,\quad z \coloneq (\norm{x_1}_2,\dots,\norm{x_N}_2) \in \R^N.
    \end{equation}
    Next, we define $[Ax]_i \coloneq ((Ax)_j)_{j=(i-1)d+1}^{id}$, $i=1,\dots,N$. For all $i$, $[Ax]_i = \sum_{j=1}^N B_{ij}x_j.$
    The norm is bounded as $\norm{[Ax]_i}_2\leq \sum_j \norm{B_{ij}x_j}_{\mathbb R^d} \leq \sum_j \norm{B_{i,j}}_{\mathbb R^d\to\mathbb R^d}\norm{x_j}_{\mathbb R^d}.$
    Therefore,
\begin{equation}
    \|Ax\|_2^2
    =
    \sum_{i=1}^N\|[Ax]_i\|_2^2 
    \le
    \sum_{i=1}^N (Cz)_i^2
    =
    \|Cz\|_2^2 
    \le
    \|C\|_{\mathbb R^{N}\to\mathbb R^{N}}^2\|z\|_2^2
    =
    \|C\|_{\mathbb R^{N}\to\mathbb R^{N}}^2\|x\|_2^2.
\end{equation}
Taking square roots and then the supremum over \(x\neq0\) proves the claim.
\end{proof}

The first Hessian mismatch comes from replacing the empirical measure
\(\rho_{\bs Z_t^{(s)}}\) by the population law \(\mu_t\) while keeping the particle locations fixed.

\begin{lemma}[Directional measure-slot error]
\label{lemma-directional-measure-slot-error}
Use the common weights $\gamma_m(D)=\mathrm L_Da_m$ from
\cref{lemma-localized-moment-smoothness-consequences}, and define
\begin{align}
 \widetilde\HessianMatrices_t^{(s)}
 &:=\HessianMatrices(\mu_t,\bs Z_t^{(s)}),\notag\\
 \Delta_{m,t}^N&:=\|M_m(\rho_{\bs Y_t})-M_m(\mu_t)\|_{\mathsf H_m},\notag\\
 \Gamma_t^N(D)&:=\sum_{m\ge1}\gamma_m(D)\Delta_{m,t}^N,\qquad
 \Lambda(D):=\sum_{m\ge1}\gamma_m(D)L_m(1+2D^{m-1}).
 \label{eq:common-moment-residual-and-lipschitz-sum}
\end{align}
On $E_{\rm pop}\cap E_{T,D}$, for every $s\in[0,1]$,
\begin{equation}
\label{eq-directional-measure-slot-bound}
 \left|\frac1N\Delta_t^\top
    (\HessianMatrices_t^{(s)}-\widetilde\HessianMatrices_t^{(s)})\Delta_t\right|
 \le\{\Lambda(D)r_t+\Gamma_t^N(D)\}r_t^2.
\end{equation}
\end{lemma}
\begin{proof}
\noindent\textbf{Step 1: bound the Hessian difference by moment displacement.}
Fix $s\in[0,1]$ and abbreviate $z_i:=Z_t^{(s),i}$, $y_i:=Y_t^i$.
Let $b_s:=\sum_m\gamma_m(D)\|M_m(\rho_{\bs Z_t^{(s)}})-M_m(\mu_t)\|$.
The two Hessian bounds in \cref{eq:local-two-hessians-moment-bound} give
\[
\begin{aligned}
 \|\nabla_W^2F(\rho_{\bs Z_t^{(s)}};z_i,z_j)
             -\nabla_W^2F(\mu_t;z_i,z_j)\|_{\mathbb R^d\to\mathbb R^d}&\le b_s/2,\\
 \|\nabla_x\nabla_WF(\rho_{\bs Z_t^{(s)}},z_i)
             -\nabla_x\nabla_WF(\mu_t,z_i)\|_{\mathbb R^d\to\mathbb R^d}&\le b_s/2.
\end{aligned}
\]
The first term of the block Hessian has the additional factor $1/N$;
the second is block diagonal. Applying
\cref{lemma-operator-norm-block-matrix} to each part therefore yields
\[
 \|\HessianMatrices_t^{(s)}-\widetilde\HessianMatrices_t^{(s)}\|_{\mathbb R^{dN}\to\mathbb R^{dN}}
 \le\left\|\frac{b_s}{2N}\mathbf1\mathbf1^\top\right\|_{\mathbb R^N\to\mathbb R^N}
      +\left\|\frac{b_s}{2}I_N\right\|_{\mathbb R^N\to\mathbb R^N}=b_s.
\]

\par\noindent\textbf{Step 2: compare the empirical moments along the interpolation.}
For the moment displacement, integrate $\fD\phi_m$ on the segment
joining $y_i$ and $z_i$, which stays in the ball of radius $D$.
The feature bound in
\cref{assumption-structured-moment-smoothness,eq:coordinate-feature-growth} gives
\[
\begin{aligned}
 \|\phi_m(z_i)-\phi_m(y_i)\|_{\mathsf H_m}
 &\le\int_0^1 L_m(1+D^{(m-1)_+})\|z_i-y_i\|_{\mathbb R^d}\,du\\
 &\le L_m(1+2D^{m-1})\|z_i-y_i\|_{\mathbb R^d}.
\end{aligned}
\]
Since $z_i-y_i=s(X_t^i-Y_t^i)$, Cauchy--Schwarz gives
\begin{equation}
\label{eq:paired-empirical-moment-displacement}
\begin{aligned}
 \|M_m(\rho_{\bs Z_t^{(s)}})-M_m(\rho_{\bs Y_t})\|_{\mathsf H_m}
 &\le\frac{L_m(1+2D^{m-1})}{N}\sum_i\|z_i-y_i\|_{\mathbb R^d}\\
 &\le L_m(1+2D^{m-1})
       \left(\frac1N\sum_i\|z_i-y_i\|_{\mathbb R^d}^2\right)^{1/2}\\
 &\le L_m(1+2D^{m-1})r_t.
\end{aligned}
\end{equation}
The triangle inequality, followed by summation with weights $\gamma_m(D)$,
now gives $b_s\le\Lambda(D)r_t+\Gamma_t^N(D)$.
Finally $N^{-1}\|\Delta_t\|_{\mathbb R^{dN}}^2=r_t^2$, so the bound on the block operator
norm proves \cref{eq-directional-measure-slot-bound}.
\end{proof}

The second Hessian mismatch comes from moving the particle locations from \(\bs Y_t\) to the interpolation \(\bs Z_t^{(s)}\) while keeping the population law fixed.

\begin{lemma}[Directional point-slot error under anti-spikiness]
\label{lemma-directional-point-slot-error}
Let
\[
    \delta_{t,i}:=X_t^i-Y_t^i,
    \qquad
    r_t^2:=\frac1N\sum_{i=1}^N\norm{\delta_{t,i}}_{\mathbb R^d}^2.
\]
Recall from \cref{lemma-directional-measure-slot-error} the
population-measure Hessian
\begin{equation}
    \widetilde\HessianMatrices_t^{(s)}
    \coloneq
    \HessianMatrices(\mu_t,\bs Z_t^{(s)}).
\end{equation}
On the event
\(E_{\rm pop}\cap E_{T,\mathrm{sp}}(A_{\mathrm{sp}})\cap E_{T,D}\),
for every \(s\in[0,1]\),
\begin{equation}
\label{eq-directional-point-slot-bound}
    \left|
    \frac1N
    \Delta_t^\top
    \bigl(
        \widetilde\HessianMatrices_t^{(s)}-
        \HessianMatrices_t^Y
    \bigr)
    \Delta_t
    \right|
    \le
    \bigl((2+A_{\mathrm{sp}})\mathrm L_D\bigr)r_t^3.
\end{equation}
\end{lemma}

\begin{proof}
Fix \(s\in[0,1]\), and write \(z_i:=Z_{t,i}^{(s)}\), \(y_i:=Y_t^i\), and
\(\delta_i:=X_t^i-Y_t^i\). Since \(z_i-y_i=s\delta_i\), we have
\(\norm{z_i-y_i}_{\mathbb R^d}\le \norm{\delta_i}_{\mathbb R^d}\).

\par\noindent\textbf{Step 1: bound the integral-Hessian contribution.}
Split the quadratic form into the integral-Hessian contribution
$Q_{\mathrm{off}}$, whose sum includes $i=j$, and the spatial-Jacobian
contribution $Q_{\mathrm{diag}}$:
\begin{align}
    Q_{\mathrm{off}}
    &:=
    \frac1{N^2}\sum_{i,j=1}^N
    \left\langle
        \delta_i,
        \bigl[
            \nabla_W^2F(\mu_t,z_i,z_j)-
            \nabla_W^2F(\mu_t,y_i,y_j)
        \bigr]
        \delta_j
    \right\rangle,\\
    Q_{\mathrm{diag}}
    &:=
    \frac1N\sum_{i=1}^N
    \left\langle
        \delta_i,
        \bigl[
            \nabla_x\nabla_WF(\mu_t,z_i)-
            \nabla_x\nabla_WF(\mu_t,y_i)
        \bigr]
        \delta_i
    \right\rangle.
\end{align}
The spatial Lipschitzness (\cref{lemma-localized-moment-smoothness-consequences}) gives
\begin{equation}
 \begin{aligned}
    \norm{
        \nabla_W^2F(\mu_t,z_i,z_j)-
        \nabla_W^2F(\mu_t,y_i,y_j)
    }_{\Loperator d}
    &\le \mathrm L_D\norm{z_i-y_i}_{\mathbb R^d}+\mathrm L_D\norm{z_j-y_j}_{\mathbb R^d}\\
    &\le \mathrm L_D\norm{\delta_i}_{\mathbb R^d}+\mathrm L_D\norm{\delta_j}_{\mathbb R^d}.
 \end{aligned}
\end{equation}
Hence
\begin{align}
    |Q_{\mathrm{off}}|
    &\le
    \frac1{N^2}\sum_{i,j=1}^N
    \norm{\delta_i}_{\mathbb R^d}\norm{\delta_j}_{\mathbb R^d}
    \bigl(\mathrm L_D\norm{\delta_i}_{\mathbb R^d}+\mathrm L_D\norm{\delta_j}_{\mathbb R^d}\bigr)\\
    &=
    \frac{\mathrm L_D}{N^2}
    \left(\sum_i\norm{\delta_i}_{\mathbb R^d}^2\right)
    \left(\sum_j\norm{\delta_j}_{\mathbb R^d}\right)
    +
    \frac{\mathrm L_D}{N^2}
    \left(\sum_i\norm{\delta_i}_{\mathbb R^d}\right)
    \left(\sum_j\norm{\delta_j}_{\mathbb R^d}^2\right).
\end{align}
By Cauchy--Schwarz,
\[
    \sum_i\norm{\delta_i}_{\mathbb R^d}
    \le
    \sqrt N\left(\sum_i\norm{\delta_i}_{\mathbb R^d}^2\right)^{1/2}
    =Nr_t,
    \qquad
    \sum_i\norm{\delta_i}_{\mathbb R^d}^2=Nr_t^2.
\]
Therefore
\begin{equation}
    |Q_{\mathrm{off}}|
    \le
    2\mathrm L_Dr_t^3.
\end{equation}
\par\noindent\textbf{Step 2: bound the spatial-Jacobian contribution.}
The spatial Lipschitzness (\cref{lemma-localized-moment-smoothness-consequences}) gives
\begin{equation}
    \norm{
        \nabla_x\nabla_WF(\mu_t,z_i)-
        \nabla_x\nabla_WF(\mu_t,y_i)
    }_{\Loperator d}
    \le
    \mathrm L_D\norm{z_i-y_i}_{\mathbb R^d}
    \le
    \mathrm L_D\norm{\delta_i}_{\mathbb R^d}.
\end{equation}
Consequently,
\begin{equation}
    |Q_{\mathrm{diag}}|
    \le
    \frac{\mathrm L_D}{N}\sum_{i=1}^N\norm{\delta_i}_{\mathbb R^d}^3
    \le
    \mathrm L_DA_{\mathrm{sp}}r_t^3
\end{equation}
on \(E_{T,\mathrm{sp}}(A_{\mathrm{sp}})\). Combining the two estimates proves
\cref{eq-directional-point-slot-bound}.
\end{proof}

Combining the measure-slot and point-slot estimates gives the directional control of the Hessian mismatch that appears in the differential inequality for \(r_t\).

\begin{corollary}[Combined directional bound for \(\iota_t^N\)]
\label{corollary-bound-iota-combined}
On $E_{\rm pop}\cap E_{T,D}\cap E_{T,\mathrm{sp}}(A_{\mathrm{sp}})$,
with the population moment event of \cref{lem:population-moment-event},
sample localization \eqref{eq:sample-localization-event}, and
anti-spikiness \eqref{eq-cubic-antispike-bound}, it holds that
\begin{equation}
\label{eq-iota-combined-bound}
    \iota_t^N
    \le
    \Bigl(\Lambda(D)+(2+A_{\mathrm{sp}})\mathrm L_D\Bigr)r_t^3
    +
    \Gamma_t^Nr_t^2.
\end{equation}
The same deterministic radius $D$ is used in the event and all coefficients.
Define $q_{4,t}:=(N^{-1}\sum_i\|X_t^i-Y_t^i\|_{\mathbb R^d}^4)^{1/4}$. In particular, if \(q_{4,t}\le A_4r_t\) for all \(t\), with $A_4\ge1$, then the same bound holds with
\(A_{\mathrm{sp}}=A_4^2\).
\end{corollary}

\begin{proof}
The directional Hessian discrepancy, recalled from
\cref{lemma-bound-r_t-given-eta-eps},
\eqref{eq:ode-hessian-discrepancy-definition}, is
$\iota_t^N=N^{-1}\sup_{s\in[0,1]}
 |\Delta_t^\top(\HessianMatrices_t^{(s)}-\HessianMatrices_t^Y)\Delta_t|$,
where $\Delta_t=\mathbf X_t-\mathbf Y_t$.
\noindent\textbf{Step 1: combine the two Hessian contributions.}
For each \(s\in[0,1]\), decompose
\begin{align}
    \HessianMatrices_t^{(s)}-
    \HessianMatrices_t^Y
    &=
    \bigl(
        \HessianMatrices_t^{(s)}-
        \widetilde\HessianMatrices_t^{(s)}
    \bigr)
    +
    \bigl(
        \widetilde\HessianMatrices_t^{(s)}-
        \HessianMatrices_t^Y
    \bigr).
\end{align}
Taking the associated quadratic form against \(\Delta_t\), dividing by \(N\), and applying the triangle inequality gives
\begin{align}
    \left|
    \frac1N\Delta_t^\top
    \bigl(\HessianMatrices_t^{(s)}-\HessianMatrices_t^Y\bigr)
    \Delta_t
    \right|
    &\le
    \left|
    \frac1N\Delta_t^\top
    \bigl(
        \HessianMatrices_t^{(s)}-
        \widetilde\HessianMatrices_t^{(s)}
    \bigr)
    \Delta_t
    \right|\\
    &\quad+
    \left|
    \frac1N\Delta_t^\top
    \bigl(
        \widetilde\HessianMatrices_t^{(s)}-
        \HessianMatrices_t^Y
    \bigr)
    \Delta_t
    \right|.
\end{align}
The first term is bounded by \cref{lemma-directional-measure-slot-error} as
\((\Lambda(D)r_t+\Gamma_t^N)r_t^2\). The second term is bounded by
\cref{lemma-directional-point-slot-error} as
\(((2+A_{\mathrm{sp}})\mathrm L_D)r_t^3\). Taking the supremum over
\(s\in[0,1]\) proves \cref{eq-iota-combined-bound}.

\par\noindent\textbf{Step 2: deduce anti-spikiness from the fourth moment.}
For $\delta_{t,i}=X_t^i-Y_t^i$, Cauchy--Schwarz gives
\[
 \frac1N\sum_i\|\delta_{t,i}\|_{\mathbb R^d}^3
 \le\left(\frac1N\sum_i\|\delta_{t,i}\|_{\mathbb R^d}^2\right)^{1/2}
     \left(\frac1N\sum_i\|\delta_{t,i}\|_{\mathbb R^d}^4\right)^{1/2}
 =r_tq_{4,t}^2\le A_4^2r_t^3.
\]
The calculation also holds for $r_t=0$, when every $\delta_{t,i}=0$.
Thus \eqref{eq-cubic-antispike-bound} holds with $A_{\mathrm{sp}}=A_4^2$.
\end{proof}

\paragraph{Bound $r_t$ by moment residual.}

The final ODE estimate packages all continuous-time residuals into the common moment fluctuation \(\Gamma_t^N\), yielding the stopped curvature-response inequality used in the main theorem.

\begin{lemma}[Common empirical moment residual and directional differential inequality for \(r_t\)]
\label{lemma-common-empirical-moment-residual-directional-rt}
Work on the population moment event of Lemma~\ref{lem:population-moment-event}, and on the events
$E_{T,\mathrm{sp}}(A_{\mathrm{sp}})$ and $E_{T,D}$, and fix an open ODE interval
\(I_\ell=(t_\ell,t_{\ell+1})\). Recall from
\cref{eq:common-moment-residual-and-lipschitz-sum} the coordinate residuals
and their weighted sum:
\begin{equation}
    \Delta_{m,t}^N
    \coloneq
    \norm{M_m(\rho_{\bs Y_t})-M_m(\mu_t)}_{\mathsf H_m},
    \qquad
    \Gamma_t^N(D)
    \coloneq
    \sum_{m\ge1}\gamma_m(D)\Delta_{m,t}^N.
\end{equation}
Here, as in \cref{eq:common-local-moment-weights},
\begin{equation}
    \gamma_m(D)
    \coloneq
    \mathrm L_Da_m.
\end{equation}
Let
\begin{equation}
\label{eq-Csp-definition}
    c_D
    \coloneq
    \Lambda(D)+(2+A_{\mathrm{sp}})\mathrm L_D.
\end{equation}
Then
\begin{equation}
    \varepsilon_t^N\le \Gamma_t^N,
\end{equation}
and, for a.e. \(t\in I_\ell\) with \(r_t>0\),
\begin{equation}
\label{eq-rt-directional-final}
    \dot r_t
    \le
    -\kappa_t^{Y,N}r_t
    +
    c_Dr_t^2
    +
    \Gamma_t^N(D)r_t
    +
    \Gamma_t^N(D).
\end{equation}
The right upper Dini derivative is
\begin{equation}
\label{eq:right-upper-dini-derivative}
 \fD^+r_t:=\limsup_{h\downarrow0}\frac{r_{t+h}-r_t}{h},
 \qquad t,t+h\in I_\ell.
\end{equation}
It agrees with $\dot r_t$ where the ordinary derivative exists.
The same upper bound holds with $\fD^+r_t$, including at $r_t=0$
as justified below. This definition is used within one ODE branch,
not across a jump.

Moreover, for the pathwise first-exit time
\begin{equation}
    \tau_R^\ell
    \coloneq
    \inf\{t\in I_\ell:r_t\ge R\},
\end{equation}
we have, for \(t<\tau_R^\ell\),
\begin{equation}
\label{eq-rt-directional-stopped}
    \dot r_t
    \le
    -\bigl(\kappa_t^{Y,N}-c_DR-\Gamma_t^N(D)\bigr)r_t
    +
    \Gamma_t^N(D).
\end{equation}
\end{lemma}

\begin{proof}
\noindent\textbf{Step 1: control the residual and positive discrepancy.}
For each \(i\), the gradient bound in \cref{eq:local-value-gradient-moment-bound} gives
\begin{equation}
\begin{aligned}
 \norm{V(\rho_{\bs Y_t},Y_t^i)-V(\mu_t,Y_t^i)}_{\mathbb R^d}
 &=\norm{\nabla_WF(\rho_{\bs Y_t},Y_t^i)-\nabla_WF(\mu_t,Y_t^i)}_{\mathbb R^d}\\
 &\le\sum_{m\ge1}\gamma_m(D)\Delta_{m,t}^N=\Gamma_t^N.
\end{aligned}
\end{equation}
Therefore
\begin{equation}
    \varepsilon_t^N
    =
    \left(
        \frac1N\sum_{i=1}^N
        \norm{V(\rho_{\bs Y_t},Y_t^i)-V(\mu_t,Y_t^i)}_{\mathbb R^d}^2
    \right)^{1/2}
    \le
    \Gamma_t^N.
\end{equation}
Next, by \cref{corollary-bound-iota-combined},
\begin{equation}
    \iota_t^N
    \le
    c_Dr_t^3+
    \Gamma_t^Nr_t^2.
\end{equation}
Before applying \cref{lemma-bound-r_t-given-eta-eps}, we verify its
differentiability premise. On $E_{\rm pop}\cap E_{T,D}$, convexity of the
radius-$D$ ball gives $Z_t^{(s),i}\in\overline B_D(0)$ for every
$s\in[0,1]$, and \cref{lemma-reachable-moment-tube} gives
$M(\rho_{\bs Z_t^{(s)}})\in\mathcal V_{K_{\rm pop},D}\subset\mathcal U$.
The $C^2$ feature map of \cref{lem:localized-weighted-feature-envelopes},
the condition $G\in C^3(\mathcal U)$ in
\cref{assumption-structured-moment-smoothness}, and the derivative formulas
\cref{eq:weighted-gradient-representation,eq:weighted-spatial-hessian-representation,eq:weighted-integral-hessian-representation}
therefore make the empirical configuration field $C^1$ along this segment.
Applying \cref{lemma-bound-r_t-given-eta-eps}, for \(r_t>0\),
\begin{align}
    \dot r_t
    &\le
    -\kappa_t^{Y,N}r_t
    +
    \frac{\iota_t^N}{r_t}
    +
    \varepsilon_t^N\\
    &\le
    -\kappa_t^{Y,N}r_t
    +
    c_Dr_t^2
    +
    \Gamma_t^Nr_t
    +
    \varepsilon_t^N\\
    &\le
    -\kappa_t^{Y,N}r_t
    +
    c_Dr_t^2
    +
    \Gamma_t^Nr_t
    +
    \Gamma_t^N.
\end{align}

\par\noindent\textbf{Step 2: handle zero discrepancy without division.}
It remains to justify the inequality at $r_t=0$, where division by $r_t$
was not used. First take $s,t$ on the same ODE branch. The pathwise bound
$\|Y_t^1-Y_s^1\|_{\mathbb R^d}\le B|t-s|$ and the first-derivative bound in
\cref{assumption-structured-moment-smoothness} give
\[
\begin{aligned}
 \|M_m(\mu_t)-M_m(\mu_s)\|_{\mathsf H_m}
 &\le\mathbb E_{\mathcal G}\|\phi_m(Y_t^1)-\phi_m(Y_s^1)\|_{\mathsf H_m}\\
 &\le BL_m|t-s|\,\mathbb E_{\mathcal G}
        (1+\|Y_t^1\|_{\mathbb R^d}^{m-1}+\|Y_s^1\|_{\mathbb R^d}^{m-1})\\
 &\le BL_m|t-s|
        [1+2\{K_{\rm pop}m^{1/\alpha}\}^{m-1}].
\end{aligned}
\]
Here $\mathbb E_{\mathcal G}$ means expectation conditional on the
population fields; the conditional moment estimate is
\eqref{eq:population-uniform-conditional-moments}. Summing with weights
$a_m$ bounds the $\mathsf H_a$ norm, since a nonnegative sequence has
$\ell^2$ norm at most its $\ell^1$ norm:
\[
 \|M(\mu_t)-M(\mu_s)\|_{\mathsf H_a}
 \le B|t-s|\sum_{m\ge1}a_mL_m
        [1+2\{K_{\rm pop}m^{1/\alpha}\}^{m-1}]<\infty.
\]
Apply \cref{lem:weighted-series-estimate},
\eqref{eq:weighted-series-sum-bound}, with $K=1+K_{\rm pop}$ and $D=1$:
$\{K_{\rm pop}m^{1/\alpha}\}^{m-1}
\le(1+K_{\rm pop})^m m^{m/\alpha}$, so the series is finite.
For either particle system $S=X$ or $S=Y$, write
$\mathbf S_t=(S_t^1,\ldots,S_t^N)$. The event $E_{T,D}$, the
branchwise $B$-Lipschitz bound, and convexity of the radius-$D$ ball give
\[
\begin{aligned}
 \|M_m(\rho_{\mathbf S_t})-M_m(\rho_{\mathbf S_s})\|_{\mathsf H_m}
 &\le\frac1N\sum_{i=1}^N
       \|\phi_m(S_t^i)-\phi_m(S_s^i)\|_{\mathsf H_m}\\
 &\le BL_m(1+D^{(m-1)_+})|t-s|.
\end{aligned}
\]
Consequently, using $D\ge1$ and the bounded-support summability in
\cref{eq:weighted-population-moment-summability},
\[
 \|M(\rho_{\mathbf S_t})-M(\rho_{\mathbf S_s})\|_{\mathsf H_a}
 \le B|t-s|\sum_{m\ge1}a_mL_m(1+D^m)<\infty.
\]
Thus the population and both empirical moment paths are continuous.
Together with continuity of $\fD G,\fD^2G$ on $\mathcal U$ from
\cref{assumption-structured-moment-smoothness}(ii), the particle drifts
in \eqref{eq:weighted-gradient-representation} are continuous on the branch.
Recall $\Delta_t=\mathbf X_t-\mathbf Y_t$ and
$V(\nu,x)=-\nabla_WF(\nu,x)$. If $r_t=0$, then $\Delta_t=0$ and
\[
 \dot\Delta_t=
 (V(\rho_{\mathbf Y_t},Y_t^i)-V(\mu_t,Y_t^i))_{i=1}^N,
 \qquad \Delta_{t+h}=h\dot\Delta_t+o(h).
\]
By the Dini definition \eqref{eq:right-upper-dini-derivative},
\[
 \fD^+r_t
 =\lim_{h\downarrow0}\frac{\|\Delta_{t+h}\|_2}{\sqrt N h}
 =\frac{\|\dot\Delta_t\|_2}{\sqrt N}
 =\varepsilon_t^N\le\Gamma_t^N.
\]
For positive $r_t$, the preceding calculation gives the same upper Dini
inequality (or its a.e. version under absolute continuity). This proves
\cref{eq-rt-directional-final}. If \(t<\tau_R^\ell\), then \(r_t<R\), so
\begin{equation}
    c_Dr_t^2
    \le
    c_DRr_t,
\end{equation}
which yields \cref{eq-rt-directional-stopped}.
\end{proof}
